\documentclass[article]{jss}

%% Requiere los ficheros de estilo de JSS (jss.cls, jsslogo.jpg y demás),
%% que se descargan de https://www.jstatsoft.org/pages/view/style
%% en el paquete jss-article-tex.zip. No se versionan aquí porque son de JSS.

\author{Luis \'Angel Barrera-Guzm\'an\\Universidad Aut\'onoma Chapingo}
\Plainauthor{Luis Angel Barrera-Guzman}

\title{\pkg{PCAPro}: Guided Principal Component Analysis in the Browser}
\Plaintitle{PCAPro: Guided Principal Component Analysis in the Browser}
\Shorttitle{\pkg{PCAPro}: Guided PCA in the Browser}

\Abstract{
  Principal component analysis is rarely limited by the eigendecomposition. What
  determines whether an analysis is defensible are the decisions surrounding it:
  how missing values are handled, whether the variables are standardised, whether
  the correlation structure justifies the analysis at all, how many components are
  retained, whether the solution is rotated, and which loadings may be interpreted.
  General-purpose implementations compute these quantities but leave the decisions
  entirely to the analyst.

  This paper introduces \pkg{PCAPro}, an implementation of the principal component
  analysis workflow written in \proglang{JavaScript} that runs entirely inside a web
  browser, without installation, without a server, and without transmitting data.
  The package implements eight retention criteria, seven factor rotations obtained
  through the gradient projection algorithm, the usual adequacy diagnostics, and the
  active/supplementary distinction of the French school of data analysis. Each
  decision point is accompanied in the interface by the methodological argument
  behind it, in Spanish and English.

  The numerical core is independent of the interface and can be used from
  \proglang{Node.js}, which makes every published result reproducible outside the
  browser. We verify the implementation against \proglang{R}
  (\code{prcomp()} and \pkg{psych}), compare its scope with existing packages, and
  illustrate its use on [PENDIENTE: caso de estudio].
}

\Keywords{principal component analysis, factor rotation, parallel analysis,
  reproducible research, statistical software, \proglang{JavaScript}}
\Plainkeywords{principal component analysis, factor rotation, parallel analysis,
  reproducible research, statistical software, JavaScript}

\Address{
  Luis \'Angel Barrera-Guzm\'an\\
  Centro Acad\'emico Regional sede Huatusco-Veracruz\\
  Universidad Aut\'onoma Chapingo\\
  Carretera Federal Huatusco-Xalapa km 6\\
  94100 Huatusco, Veracruz, Mexico\\
  E-mail: \email{lbarrerag@chapingo.mx}\\
  ORCID: \url{https://orcid.org/0000-0001-8057-2583}
}

\begin{document}

%% -------------------------------------------------------------------------
\section{Introduction} \label{sec:intro}

Principal component analysis \citep{pearson1901, hotelling1933, jolliffe2002} is
among the most widely used multivariate techniques in the biological and
agricultural sciences, and it is often the first multivariate method a student
encounters. It is also among the most frequently misapplied.

The computation itself is not the difficulty. Every general-purpose statistical
environment provides a reliable eigendecomposition, and the numerical problem has
been settled for decades. What decides whether a published principal component
analysis is defensible is a sequence of choices made before and after that
computation: how missing values are treated, whether variables are standardised,
whether the correlation structure justifies the analysis at all, how many
components are retained, whether the solution is rotated and by which method, and
which loadings are large enough to interpret. Each of these changes the result,
and none of them is decided by a function call.

Existing implementations compute the quantities needed to make those decisions and
then stop. \pkg{FactoMineR} \citep{le2008} reports the eigenvalues; deciding how many
to keep is left to the analyst. \pkg{psych} \citep{psych} reports the
Kaiser--Meyer--Olkin index; deciding whether $0.62$ is acceptable is left to the
analyst. This division of labour is entirely reasonable for a statistician. It is
where non-specialist users go wrong.

\pkg{PCAPro} takes a different position: the methodological argument belongs next to
the number it justifies. The package reports eight retention criteria side by side
rather than a single answer, because the criteria disagree and the disagreement is
informative; it states which assumption is at risk and what to do about it rather
than printing a diagnostic and moving on; and it declines to present a decision as
more certain than it is.

Two further constraints shaped the implementation and are not incidental to it.
First, the software runs with no installation: it is a directory of static files
served by any web server or opened directly from disk. This is a functional
requirement rather than a convenience, because a substantial part of the intended
audience works in shared university computer laboratories where nothing can be
installed. Second, the interface, the embedded methodological guidance and the
generated report exist in full in Spanish as well as English, because teaching
material for multivariate statistics in Spanish is scarce and the language barrier
compounds the tooling barrier.

The rest of this paper is organised as follows.
Section~\ref{sec:methods} reviews the decision points of the principal component
analysis workflow and the criteria \pkg{PCAPro} implements at each one.
Section~\ref{sec:implementation} describes the design of the software and the
trade-offs behind it.
Section~\ref{sec:comparison} compares \pkg{PCAPro} with existing implementations and
verifies its numerical output against \proglang{R}.
Section~\ref{sec:illustration} illustrates its use on a complete analysis.
Section~\ref{sec:summary} concludes.

%% -------------------------------------------------------------------------
\section{The decision points of a principal component analysis} \label{sec:methods}

%% PENDIENTE: esta seccion recorre las decisiones y los criterios implementados.
%% El material ya esta escrito en las 34 secciones teoricas de la aplicacion
%% (index.html) y hay que trasladarlo aqui, en registro de articulo y con las
%% citas de paper/jss/ref.bib. Subsecciones previstas:

\subsection{Adequacy of the data}
%% Bartlett \citep{bartlett1950}, KMO y MSA \citep{kaiser1974}, determinante,
%% multicolinealidad, atipicos de Mahalanobis, razon n:p.

\subsection{Standardisation and the covariance/correlation choice}
%% Las siete opciones de escalado y cuando aplica cada una.

\subsection{How many components to retain}
%% Los ocho criterios: analisis paralelo de Horn \citep{horn1965} con la regla del
%% percentil 95 \citep{glorfeld1995}, Kaiser--Guttman, Jolliffe, baston roto
%% \citep{frontier1976}, MAP de Velicer \citep{velicer1976}, codo del scree
%% \citep{cattell1966} y umbrales de varianza acumulada. Discutir por que se
%% presentan juntos y como se pondera el consenso.

\subsection{Rotation}
%% Familia ortomax y rotaciones oblicuas \citep{kaiser1958, hendrickson1964},
%% obtenidas por el algoritmo de proyeccion de gradiente
%% \citep{jennrich2001, jennrich2002}. Matrices de patron, estructura y Phi.

\subsection{Interpretation}
%% cos2, contribuciones, elementos suplementarios, valores test.

%% -------------------------------------------------------------------------
\section{Implementation} \label{sec:implementation}

\subsection{Architecture}

\pkg{PCAPro} is a directory of static files. It uses no framework, no module
bundler and no build step, and every component is attached to the global object so
that the application also runs from a \code{file://} URL. This is a deliberate
trade-off. The cost is real: there is no module system, no tree shaking, and the
interface is constructed by explicit DOM manipulation. The benefit is that the
software has no toolchain that can rot, can be served from a memory stick or an
intranet, and can be read file by file by a reviewer or a student without
installing anything. For software whose purpose is to work where nothing can be
installed, freedom from dependencies is a functional requirement rather than an
aesthetic preference.

The only third-party component is \pkg{SheetJS}, used exclusively to read
\code{.xlsx} workbooks; it can be vendored locally for a fully offline
installation.

\subsection{The numerical core}

The numerical engine was written rather than wrapped. No \proglang{JavaScript}
library offered the combination of retention criteria and rotation algorithms
required, and a self-contained implementation makes every step auditable in the
same language as the interface.

Three files constitute the core and are independent of the interface: linear
algebra and distributions, the rotation algorithms, and the factor coordinates
with their quality measures. They export themselves both as browser scripts and as
\proglang{Node.js} modules, so there is a single implementation with two ways of
being loaded rather than two implementations that can drift apart:

\begin{Code}
const { S, Rot, Fac } = require("./lib");

const R = S.corrMatrix(columns);
const eig = S.eigenSym(R);
const kmo = S.kmo(R);
const rot = Rot.rotate(loadings, "varimax", { normalize: true });
\end{Code}

This is what makes the results of the present paper reproducible without a browser;
Section~\ref{sec:comparison} relies on it.

The eigendecomposition uses cyclic Jacobi rotations on the symmetric dispersion
matrix. The factor rotations use the gradient projection algorithm of
\citet{jennrich2001, jennrich2002}, which handles the orthogonal and oblique cases
within a single framework and admits the orthomax family, quartimin and oblimin
through their respective criterion functions; \pkg{promax}
\citep{hendrickson1964} is obtained from a varimax solution followed by a
least-squares fit to the powered target.

\subsection{Graphics}

Figures are drawn by an in-house SVG generator rather than a charting library.
Charting libraries optimise for rendering; publication figures require editability
and vector export at arbitrary resolution, with the exported file corresponding
exactly to what was edited on screen. Owning the SVG generation makes that
correspondence exact and allows every figure to carry the same editing controls:
titles, axis labels, palettes, themes, ellipse type, typeface and four independent
font-size controls.

\subsection{Verification}

Correctness is checked by an automated suite of 60 tests. Every expectation is
either a reference value published by an independent implementation or an
algebraic invariant that must hold regardless of the data, such as the eigenvalues
summing to the trace or communalities being preserved by an orthogonal rotation.
No expectation was recorded from a previous run of \pkg{PCAPro} itself: comparing
software against its own past output establishes stability, not correctness. The
varimax implementation is additionally checked against an exhaustive sweep of the
rotation angle, confirming that the gradient projection algorithm reaches the
global maximum of its own criterion rather than a local one. The suite runs
automatically in a headless browser on every change to the repository.

%% -------------------------------------------------------------------------
\section{Comparison with existing implementations} \label{sec:comparison}

%% PENDIENTE: JSS pide explicitamente que la comparacion se apoye en
%% "empirical illustrations". Hay que ampliar esta seccion con una tabla que
%% compare, sobre los mismos datos, los valores que devuelven PCAPro, prcomp,
%% FactoMineR y psych. El guion lib/reproduce-iris.js ya produce la mitad.

\subsection{Scope}

%% Trasladar aqui la comparacion ya escrita: FactoMineR + factoextra, psych,
%% scikit-learn, JASP, jamovi, PAST y los paquetes comerciales; con la
%% justificacion de construir en lugar de contribuir.

\subsection{Numerical agreement}

%% Tabla con los valores de referencia y los obtenidos, generada por
%% lib/reproduce-iris.js.

%% -------------------------------------------------------------------------
\section{Illustration} \label{sec:illustration}

%% PENDIENTE, y es lo que decide el articulo. JSS exige "an enlightening
%% non-trivial case study". Iris no sirve: es trivial y todo el mundo lo usa.
%% Hace falta un conjunto de datos real, del autor, con suficientes variables
%% cuantitativas y una estructura que merezca interpretarse; y hay que poder
%% distribuirlo con el material de replicacion.

%% -------------------------------------------------------------------------
\section{Summary and discussion} \label{sec:summary}

%% PENDIENTE.

%% -------------------------------------------------------------------------
\section*{Computational details}

The results in this paper were obtained using \pkg{PCAPro}~1.0.1
\citep{pcapro}. Reference values were computed in \proglang{R} using
\code{prcomp()} from \pkg{stats} \citep{R} and \pkg{psych} \citep{psych}.
The replication script \code{lib/reproduce-iris.js} reproduces every numerical
result reported here and is included with the software.

\bibliography{ref}

\end{document}
